\section{Applications of Fourier series}

    \subsection{Fourier expansion of square wave}
        A square wave may be represented mathematically by the following piece-wise continuous function:
        \begin{align*}
            f(x)=\begin{cases}
                {-k} & {\text{for }-\pi<x<0} \\
                {k} & {\text{for }0<x<\pi}
            \end{cases}
        \end{align*}
        and $f(x+2\pi)=f(x)$.

        The square wave function is plotted below:
        \begin{figure}[h!]
            \centering
            \includegraphics[width=0.55\textwidth]{figures/fig_square-wave.png}
        \end{figure}

        In order to obtain the Fourier series expansion of the square wave function, we need to first check if the Dirichlet's conditions are satisfied.
        \begin{itemize}
            \item $f(x)$ has a finite number of maxima and minima in the range $-\pi<0<\pi$.
            \item $f(x)$ has a finite number of discontinuities in the range $-\pi<0<\pi$.
            \item $f(x)$ is finite and is therefore absolutely integrable over the range $-\pi<0<\pi$.
        \end{itemize}
        Because $f(x)$ satisfies the Dirichlet's conditions, we know that it can be expanded as a Fourier series. Now we can find its Fourier coefficients.

        \begin{itemize}
            \item To find $a_0$:
            \begin{align*}
                a_0&=\dfrac{1}{2\pi}\int\limits_{-\pi}^{\pi}{f(x)\,\dd{x}} \\
                    &=\dfrac{1}{2\pi}\qty[\int\limits_{-\pi}^{0}{(-k)\,\dd{x}}+\int\limits_{0}^{\pi}{(k)\,\dd{x}}]=\dfrac{1}{2\pi}\qty[-k\int\limits_{-\pi}^{0}{\dd{x}}+k\int\limits_{0}^{\pi}{\dd{x}}]\\
                    &=\dfrac{1}{2\pi}\qty[-k\qty{0-(-\pi)}+k(\pi-0)]=\dfrac{1}{2\pi}\qty[-k\pi+k\pi] \\
                \implies a_0&=0
            \end{align*}

            \item To find $a_n$:
            \begin{align*}
                a_n&=\dfrac{1}{\pi}\int\limits_{-\pi}^{\pi}{f(x)\cos{nx}\,\dd{x}} \\
                    &=\dfrac{1}{\pi}\qty[\int\limits_{-\pi}^{0}{(-k)\cos{nx}\,\dd{x}}+\int\limits_{0}^{\pi}{(k)\cos{nx}\,\dd{x}}] \\
                    &=\dfrac{k}{\pi}\qty[-\int\limits_{-\pi}^{0}{\cos{nx}\,\dd{x}}+\int\limits_{0}^{\pi}{\cos{nx}\,\dd{x}}] \\
                    &=\dfrac{k}{\pi}\qty[-\left.\dfrac{\sin{nx}}{n}\right|_{-\pi}^{0}+\left.\dfrac{\sin{nx}}{n}\right|_{0}^{\pi}] \\
                    &=\dfrac{k}{\pi}\qty[\dfrac{\cancelto{0}{\sin{(-n\pi)}}-\cancelto{0}{\sin{0}}}{n}+\dfrac{\cancelto{0}{\sin{(n\pi)}}-\cancelto{0}{\sin{0}}}{n}] \\
                \implies a_n&=0\qq{(for all $n$)}
            \end{align*}

            \item To find $b_n$:
            \begin{align*}
                b_n&=\dfrac{1}{\pi}\int\limits_{-\pi}^{\pi}{f(x)\sin{nx}\,\dd{x}} \\
                    &=\dfrac{1}{\pi}\qty[\int\limits_{-\pi}^{0}{(-k)\sin{nx}\,\dd{x}}+\int\limits_{0}^{\pi}{(k)\sin{nx}\,\dd{x}}] \\
                    &=\dfrac{k}{\pi}\qty[-\int\limits_{-\pi}^{0}{\sin{nx}\,\dd{x}}+\int\limits_{0}^{\pi}{\sin{nx}\,\dd{x}}] \\
                    &=\dfrac{k}{\pi}\qty[\left.\dfrac{\cos{nx}}{n}\right|_{-\pi}^{0}-\left.\dfrac{\cos{nx}}{n}\right|_{0}^{\pi}] \\
                    &=\dfrac{k}{\pi}\qty[\dfrac{{\cos{0}}-{\cos{(-n\pi)}}}{n}-\dfrac{{\cos{(n\pi)}}-{\cos{0}}}{n}] \\
                    &=\dfrac{k}{n\pi}\qty(\cos{0}-\cos{n\pi}-\cos{n\pi}+\cos{0}) \\
                \implies b_n&=\dfrac{2k}{n\pi}\qty(1-\cos{n\pi})
            \end{align*}
            Now, we know that
            \begin{align*}
                \cos{n\pi}&=\begin{cases}
                    {-1;} & {\text{for odd }n} \\
                    {1;} & {\text{for even }n}
                \end{cases} \\
                \implies 1-\cos{n\pi}&=\begin{cases}
                    {2;} & {\text{for odd }n} \\
                    {0;} & {\text{for even }n}
                \end{cases}
            \end{align*}
            Hence, we obtain
            \begin{align*}
                b_n&=\begin{cases}
                    {\dfrac{4k}{n\pi};} & {\text{for odd }n} \\
                    {0;} & {\text{for even }n}
                \end{cases}
            \end{align*}
        \end{itemize}
        This gives us
            \begin{align*}
                a_0&=0,\quad a_n=0\text{ (for all $n$)} \\
                b_1&=\dfrac{4k}{\pi},\quad b_2=0,\quad b_3=\dfrac{4k}{3\pi},\quad b_4=0,\quad b_5=\dfrac{4k}{5\pi},\quad b_6=0,\quad\cdots
            \end{align*}

        Therefore, the Fourier series expansion of the square wave function is
        \begin{align*}
            f(x)&=b_1\sin{x}+b_3\sin{3x}+b_5\sin{5x}+\cdots \\
                &=\dfrac{4k}{\pi}\sin{x}+\dfrac{4k}{3\pi}\sin{3x}+\dfrac{4k}{5\pi}\sin{5x}+\cdots \\
                &=\dfrac{4k}{\pi}\qty(\sin{x}+\dfrac{\sin{3x}}{3}+\dfrac{\sin{5x}}{5}+\cdots) \\
            \implies\Aboxed{f(x)&=\dfrac{4k}{\pi}\sum_{m=0}^{\infty}{\dfrac{\sin{(2m+1)x}}{2m+1}}}
        \end{align*}

    \setcounter{equation}{0}
    \subsection{Fourier expansion of triangular wave}
         A triangular wave may be represented mathematically by the following piece-wise continuous function:
        \begin{align*}
            f(x)=\begin{cases}
                {\dfrac{k}{\pi}x} & {\text{for }0<x<\pi} \\
                {} & {} \\
                {\dfrac{k}{\pi}(2\pi-x)} & {\text{for }\pi<x<2\pi}
            \end{cases}
        \end{align*}
        and $f(x+2\pi)=f(x)$.

        The square wave function is plotted below:
        \begin{figure}[h!]
            \centering
            \includegraphics[width=0.55\textwidth]{figures/fig_triangular-wave.png}
        \end{figure}

        We can re-write the triangular function over the range $-\pi<x<\pi$ as follows:
        \begin{align*}
            f(x)=\begin{cases}
                {-\dfrac{k}{\pi}x} & {\text{for }-\pi<x<0} \\
                {} & {} \\
                {\dfrac{k}{\pi}x} & {\text{for }0<x<\pi}
            \end{cases}
        \end{align*}
        and $f(x+2\pi)=f(x)$.

        In order to obtain the Fourier series expansion of the square wave function, we need to first check if the Dirichlet's conditions are satisfied.
        \begin{itemize}
            \item $f(x)$ has a finite number of maxima and minima in the range $-\pi<0<\pi$.
            \item $f(x)$ has a finite number of discontinuities in the range $-\pi<0<\pi$.
            \item $f(x)$ is finite and is therefore absolutely integrable over the range $-\pi<0<\pi$.
        \end{itemize}
        Because $f(x)$ satisfies the Dirichlet's conditions, we know that it can be expanded as a Fourier series. Now we can find its Fourier coefficients.

        \begin{itemize}
            \item To find $a_0$:
            \begin{align*}
                a_0&=\dfrac{1}{2\pi}\int\limits_{-\pi}^{\pi}{f(x)\,\dd{x}} \\
                    &=\dfrac{1}{2\pi}\qty[\int\limits_{-\pi}^{0}{\qty(-\dfrac{k}{\pi}x)\,\dd{x}}+\int\limits_{0}^{\pi}{\qty(\dfrac{k}{\pi}x)\,\dd{x}}]=\dfrac{1}{2\pi}\qty[-\dfrac{k}{\pi}\int\limits_{-\pi}^{0}{x\dd{x}}+\dfrac{k}{\pi}\int\limits_{0}^{\pi}{x\dd{x}}]\\
                    &=\dfrac{k}{2\pi^2}\qty[\left.-\dfrac{x^2}{2}\right|_{-\pi}^{0}+\left.\dfrac{x^2}{2}\right|_{0}^{\pi}]=\dfrac{k}{4\pi^2}\qty[\pi^2+\pi^2]=\dfrac{k}{4\pi^2}\cdot 2\pi^2 \\
                \implies a_0&=\dfrac{k}{2}
            \end{align*}

            \item To find $a_n$:
            \begin{align}
                a_n&=\dfrac{1}{\pi}\int\limits_{-\pi}^{\pi}{f(x)\cos{nx}\,\dd{x}} \nonumber\\
                    &=\dfrac{1}{\pi}\qty[\int\limits_{-\pi}^{0}{\qty(-\dfrac{k}{\pi}x)\cos{nx}\,\dd{x}}+\int\limits_{0}^{\pi}{\qty(\dfrac{k}{\pi}x)\cos{nx}\,\dd{x}}] \nonumber\\
                \implies a_n&=\dfrac{k}{\pi^2}\qty[-\int\limits_{-\pi}^{0}{x\cos{nx}\,\dd{x}}+\int\limits_{0}^{\pi}{x\cos{nx}\,\dd{x}}] \label{eqn:triangular:an}
            \end{align}
            From integration-by-parts, we have
            \begin{align*}
                \int{u\dd{v}}&=uv-\int{v\dd{u}}
            \end{align*}
            We choose
            \begin{align*}
                u=x&\quad\implies\dd{u}=\dd{x} \\
                \dd{v}=\cos{nx}\dd{x}&\quad\implies v=\dfrac{1}{n}\sin{nx}
            \end{align*}
            Substituting above, we get
            \begin{align*}
                \int{x\cos{nx}\dd{x}}&=\dfrac{x}{n}\sin{nx}-\dfrac{1}{n}\int{\sin{nx}\dd{x}} \\
                \implies\int{x\cos{nx}\dd{x}}&=\dfrac{x}{n}\sin{nx}+\dfrac{1}{n^2}\cos{nx}
            \end{align*}
            Then the definite integrals become:\vspace{-3ex}
            \begin{multicols}{2}

                \begin{align*}
                    &\int\limits_{-\pi}^{0}{x\cos{nx}\dd{x}}=\left.\dfrac{x}{n}\sin{nx}\right|_{-\pi}^{0}+\left.\dfrac{1}{n^2}\cos{nx}\right|_{-\pi}^{0} \\
                    &=\dfrac{\qty[0-(-\pi)\sin n(-\pi)]}{n}+\dfrac{\qty[\cos{0}-\cos{n(-\pi)}]}{n^2} \\
                    &=-\dfrac{\pi}{n}\underbrace{\sin{n\pi}}_{=0}+\dfrac{1-\cos{n\pi}}{n^2} \\
                    &\implies\int\limits_{-\pi}^{0}{x\cos{nx}\dd{x}}=\dfrac{1-\cos{n\pi}}{n^2}
                \end{align*}

                \begin{align*}
                    &\int\limits_{0}^{\pi}{x\cos{nx}\dd{x}}=\left.\dfrac{x}{n}\sin{nx}\right|_{0}^{\pi}+\left.\dfrac{1}{n^2}\cos{nx}\right|_{0}^{\pi} \\
                    &=\dfrac{\qty[\pi\sin n\pi]-0}{n}+\dfrac{\qty[\cos{n\pi}-\cos{0}]}{n^2} \\
                    &=\dfrac{\pi}{n}\underbrace{\sin{n\pi}}_{=0}+\dfrac{\cos{n\pi}-1}{n^2} \\
                    &\implies\int\limits_{0}^{\pi}{x\cos{nx}\dd{x}}=\dfrac{\cos{n\pi}-1}{n^2}
                \end{align*}
                
            \end{multicols}

            Substituting in equation (\ref{eqn:triangular:an}), we get
            \begin{align*}
                a_n&=\dfrac{k}{\pi^2}\qty[-\qty(\dfrac{1-\cos{n\pi}}{n^2})+\qty(\dfrac{\cos{n\pi}-1}{n^2})] \\
                \implies a_n&=\dfrac{2k}{n^2\pi^2}\qty(\cos{n\pi}-1)
            \end{align*}
            
            Now, we know that
            \begin{align*}
                \cos{n\pi}&=\begin{cases}
                    {-1;} & {\text{for odd }n} \\
                    {1;} & {\text{for even }n}
                \end{cases} \\
                \implies \cos{n\pi}-1&=\begin{cases}
                    {-2;} & {\text{for odd }n} \\
                    {0;} & {\text{for even }n}
                \end{cases}
            \end{align*}

            Hence, we obtain
            \begin{align*}
                a_n&=\begin{cases}
                    {-\dfrac{4k}{n^2\pi^2};} & {\text{for odd }n} \\
                    {0;} & {\text{for even }n}
                \end{cases}
            \end{align*}

            This gives us
            \begin{align*}
                a_1&=-\dfrac{4k}{\pi^2},\quad a_2=0,\quad a_3=-\dfrac{4k}{9\pi^2},\quad a_4=0,\quad a_5=-\dfrac{4k}{25\pi^2},\quad a_6=0,\quad\cdots
            \end{align*}

            \item To find $b_n$:
            \begin{align}
                b_n&=\dfrac{1}{\pi}\int\limits_{-\pi}^{\pi}{f(x)\sin{nx}\,\dd{x}} \nonumber\\
                    &=\dfrac{1}{\pi}\qty[\int\limits_{-\pi}^{0}{\qty(-\dfrac{k}{\pi}x)\sin{nx}\,\dd{x}}+\int\limits_{0}^{\pi}{\qty(\dfrac{k}{\pi}x)\sin{nx}\,\dd{x}}] \nonumber\\
                \implies b_n&=\dfrac{k}{\pi^2}\qty[-\int\limits_{-\pi}^{0}{x\sin{nx}\,\dd{x}}+\int\limits_{0}^{\pi}{x\sin{nx}\,\dd{x}}] \label{eqn:triangular:bn}
            \end{align}

            Again from integration-by-parts, we have
            \begin{align*}
                \int{u\dd{v}}&=uv-\int{v\dd{u}}
            \end{align*}
            We choose
            \begin{align*}
                u=x&\quad\implies\dd{u}=\dd{x} \\
                \dd{v}=\sin{nx}\dd{x}&\quad\implies v=-\dfrac{1}{n}\cos{nx}
            \end{align*}

            Substituting above, we get
            \begin{align*}
                \int{x\sin{nx}\dd{x}}&=-\dfrac{x}{n}\cos{nx}+\dfrac{1}{n}\int{\cos{nx}\dd{x}} \\
                \implies\int{x\sin{nx}\dd{x}}&=-\dfrac{x}{n}\cos{nx}+\dfrac{1}{n^2}\sin{nx}
            \end{align*}

            Then the definite integrals become:\vspace{-3ex}
            \begin{multicols}{2}

                \begin{align*}
                    &\int\limits_{-\pi}^{0}{x\sin{nx}\dd{x}}=\left.-\dfrac{x}{n}\cos{nx}\right|_{-\pi}^{0}+\left.\dfrac{1}{n^2}\sin{nx}\right|_{-\pi}^{0} \\
                    &=\dfrac{\qty[(-\pi)\cos{n(-\pi)-0}]}{n}+\dfrac{\qty[\sin{0}-\sin{n(-\pi)}]}{n^2} \\
                    &=-\dfrac{\pi}{n}\cos{n\pi}+\dfrac{1}{n^2}\underbrace{\sin{n\pi}}_{=0} \\
                    &\implies\int\limits_{-\pi}^{0}{x\sin{nx}\dd{x}}=-\dfrac{\pi}{n}\cos{n\pi}
                \end{align*}

                \begin{align*}
                    &\int\limits_{0}^{\pi}{x\sin{nx}\dd{x}}=\left.-\dfrac{x}{n}\cos{nx}\right|_{0}^{\pi}+\left.\dfrac{1}{n^2}\sin{nx}\right|_{0}^{\pi} \\
                    &=\dfrac{\qty[0-\pi\cos{n\pi}]}{n}+\dfrac{\qty[\sin{n\pi}-\sin{0}]}{n^2} \\
                    &=-\dfrac{\pi}{n}\cos{n\pi} \\
                    &\implies\int\limits_{0}^{\pi}{x\sin{nx}\dd{x}}=-\dfrac{\pi}{n}\cos{n\pi}
                \end{align*}
                
            \end{multicols}

            Substituting in equation (\ref{eqn:triangular:bn}), we get
            \begin{align*}
                b_n&=\dfrac{k}{\pi^2}\qty[-\qty(-\dfrac{\pi}{n}\cos{n\pi})+\qty(-\dfrac{\pi}{n}\cos{n\pi})] \\
                    &=\dfrac{k}{n\pi}\qty(\cos{n\pi}-\cos{n\pi}) \\
                \implies b_n&=0\qq{(for all $n$)}
            \end{align*}
        \end{itemize}

        Hence, the Fourier coefficients are as follows:
        \begin{align*}
            a_0&=\dfrac{k}{2} \\
            a_1&=-\dfrac{4k}{\pi^2},\quad a_2=0,\quad a_3=-\dfrac{4k}{9\pi^2},\quad a_4=0,\quad a_5=-\dfrac{4k}{25\pi^2},\quad a_6=0,\quad\cdots \\
            b_n&=0\qq{(for all $n$)}
        \end{align*}

        Therefore, the Fourier series expansion of the triangular wave function is
        \begin{align*}
            f(x)&=a_0+a_1\cos{x}+a_3\cos{3x}+a_5\cos{5x}+\cdots \\
                &=\dfrac{k}{2}-\dfrac{4k}{\pi^2}\cos{x}-\dfrac{4k}{9\pi^2}\cos{3x}-\dfrac{4k}{25\pi^2}\cos{5x}+\cdots \\
                &=\dfrac{k}{2}-\dfrac{4k}{\pi^2}\qty(\cos{x}+\dfrac{\cos{3x}}{3^2}+\dfrac{\cos{5x}}{5^2}+\cdots) \\
            \implies\Aboxed{f(x)&=\dfrac{k}{2}-\dfrac{4k}{\pi^2}\sum_{m=0}^{\infty}{\dfrac{\cos{(2m+1)x}}{(2m+1)^2}}}
        \end{align*}

    % \newpage
    % For an asymmetric square well potential, given a particle with the wave-function
    % \begin{align*}
    %     \Psi(x,0)=\sqrt{\dfrac{1}{3L}}\sin\qty(\dfrac{\pi x}{L})+C\sin\qty(\dfrac{3\pi x}{L})
    % \end{align*}
    % where $L$ is the width of the well.
    % \begin{enumerate}
    %     \item Find $C$ so that $\Psi(x,0)$ represents a normalized state.
    %     \item Obtain $\Psi(x,t)$ and determine whether it represents a stationary or non-stationary state.
    %     \item Find the average energy $\varepsilon_\text{av}$ of the particle at time $t$ and state whether $\varepsilon_\text{av}$ depends on time, given that the ground state energy of the particle is 1 MeV.
    % \end{enumerate}

    % \newpage
    % The wave function of a free particle of mass $m$ which is constrained to move in the region\\ $-L<x<L$ is given by
    % $$\psi(x)=A(L+x)(L-x)$$
    % where $A$ is the normalization constant. Find the probability that the particle will have\\ energy $\dfrac{\pi^2\hbar^2}{2mL^2}$.

    % \newpage
    % An EM wave (of wavelength $\lambda_0$ in free space) travels through an absorbing medium with dielectric permittivity $\epsilon=\epsilon_R+i\epsilon_I$, where $\dfrac{\epsilon_I}{\epsilon_R}=\sqrt{3}$. If the skin depth is $\dfrac{\lambda_0}{4\pi}$, calculate the ratio of amplitude of the electric field to that of the magnetic field in the medium.